\subsection{希尔伯特--施密特积分算子}

在本节中，X 与 Y 是局部紧拓扑空间。设 \(\mu\) 是 X 上的正测度，\(\nu\) 是 Y 上的正测度。我们分别用 \(\mathrm{L}^{2}(\mathrm{X})\) 、\(\mathrm{L}^{2}(\mathrm{Y})\) 与 \(\mathrm{L}^{2}(\mathrm{X} \times \mathrm{Y})\) 表示空间 \(\mathrm{L}^{2}(\mathrm{X}, \mu)\) 、\(\mathrm{L}^{2}(\mathrm{Y}, \nu)\) 与 \(\mathrm{L}^{2}(\mathrm{X} \times \mathrm{Y}, \mu \otimes \nu)\) ，对 \(\mathcal{L}^{2}(\mathrm{X})\) 、\(\mathcal{L}^{2}(\mathrm{Y})\) 与 \(\mathcal{L}^{2}(\mathrm{X} \times \mathrm{Y})\) 亦同样处理。

我们回顾（参见 III, p. 26 的 no. 4），对每个 \(N \in \mathcal{L}^{2}(X \times Y)\) ，存在唯一的连续线性映射 \(u_{\mathrm{N}}\) （从 \(L^{2}(X)\) 到 \(L^{2}(Y)\) 中的连续线性映射），使得

\[
\langle g \mid u _ {\mathrm{N}} (f) \rangle = \int_ {\mathrm{X} \times \mathrm{Y}} \overline {{g (y)}} \mathrm{N} (x, y) f (x) d (\mu \otimes \nu) (x, y)\tag{12}
\]

对所有 \(f \in \mathrm{L}^{2}(\mathrm{X})\) 与每个 \(g \in \mathrm{L}^{2}(\mathrm{Y})\) 成立。映射 \(u_{N}\) 是紧的（III, p. 33 的推论 1）。映射 \(N \mapsto u_{N}\) 过渡到商，诱导出从 \(\mathrm{L}^{2}(\mathrm{X} \times \mathrm{Y})\) 到 \(\mathcal{L}(\mathrm{L}^{2}(\mathrm{X});\mathrm{L}^{2}(\mathrm{Y}))\) 中的连续单射线性映射（III, p. 30 的命题 5）。

我们用 \(\theta\) 表示从 \(\mathcal{L}^2 (\mathrm{X})\otimes \mathcal{L}^2 (\mathrm{Y})\) 到 \(\mathcal{L}^2 (\mathrm{X}\times \mathrm{Y})\) 中的唯一线性映射，它对所有 \(f\in \mathcal{L}^2 (\mathrm{X})\) 与 \(g\in \mathcal{L}^2 (\mathrm{Y})\) 把 \(f\otimes g\) 映为由 \((x,y)\mapsto f(x)g(y)\) 定义的函数。

引理 10. --- 映射 \(\theta\) 过渡到商，定义了线性映射 \(\widetilde{\theta}\) （从 \(\mathrm{L}^2 (\mathrm{X})\otimes \mathrm{L}^2 (\mathrm{Y})\) 到 \(\mathrm{L}^2 (\mathrm{X}\times \mathrm{Y})\) 中的线性映射），且存在唯一的从 \(\mathrm{L}^2 (\mathrm{X})\widehat{\otimes}_2\mathrm{L}^2 (\mathrm{Y})\) 到 \(\mathrm{L}^2 (\mathrm{X}\times \mathrm{Y})\) 上的等距同构，它在 \(\mathrm{L}^2 (\mathrm{X})\otimes \mathrm{L}^2 (\mathrm{Y})\) 上与之相合。

第一个断言是初等的。我们来证明第二个。

设 \((f_i)_{i\in \mathrm{I}}\) 与 \((g_j)_{j\in \mathrm{J}}\) 分别是 \(\mathrm{L}^2 (\mathrm{X})\) 与 \(\mathrm{L}^2 (\mathrm{Y})\) 的标准正交基。族 \((\overline{f}_i\otimes g_j)_{(i,j)\in \mathrm{I}\times \mathrm{J}}\) 是 \(\mathrm{L}^2 (\mathrm{X})\widehat{\otimes}_2\mathrm{L}^2 (\mathrm{Y})\) 的标准正交基（EVT, V, p. 29，推论 1），而族 \((\widetilde{\theta} (\overline{f}_i\otimes g_j))_{(i,j)\in \mathrm{I}\times \mathrm{J}}\) 在 \(\mathrm{L}^2 (\mathrm{X}\times \mathrm{Y})\) 中标准正交。因此映射 \(\widetilde{\theta}\) 由连续性可扩张为从 \(\mathrm{L}^2 (\mathrm{X})\widehat{\otimes}_2\mathrm{L}^2 (\mathrm{Y})\) 到 \(\mathrm{L}^2 (\mathrm{X}\times \mathrm{Y})\) 中的等距线性映射。剩下要证明这个扩张是满射。

为此，只需证明 \(\widetilde{\theta}\) 的像在 \(\mathrm{L}^{2}(\mathrm{X}\times\mathrm{Y})\) 中稠密。设 N 是 \(\mathrm{L}^{2}(\mathrm{X}\times\mathrm{Y})\) 中与 \(\widetilde{\theta}\) 的像正交的元素。对所有 \(i\in I\) 与 \(j\in J\) ，我们有 \(\langle g_{j}\mid u_{\mathrm{N}}(f_{i})\rangle=\langle\widetilde{\theta}(\overline{f}_{i}\otimes g_{j})\mid\mathrm{N}\rangle=0\)。由公式 (12) 可得 \(u_{N}=0\) ，从而 N=0。

前一引理证实了 EVT, V, p. 29 例 2 中所述的断言。

在下文中，我们将借助引理 10 的等距同构把 \(\mathrm{L}^{2}(\mathrm{X})\widehat{\otimes}_{2}\mathrm{L}^{2}(\mathrm{Y})\) 与 \(\mathrm{L}^{2}(\mathrm{X}\times\mathrm{Y})\) 等同起来。

命题 19. --- 映射 \(N\mapsto u_{\mathrm{N}}\) 是从 \(\mathrm{L}^{2}(\mathrm{X}\times\mathrm{Y})\) 到空间 \(\mathcal{L}_{2}(\mathrm{L}^{2}(\mathrm{X});\mathrm{L}^{2}(\mathrm{Y}))\) （从 \(\mathrm{L}^{2}(\mathrm{X})\) 到 \(\mathrm{L}^{2}(\mathrm{Y})\) 中的希尔伯特--施密特映射所成的空间）上的等距同构。

从 \(\mathrm{L}^{2}(\mathrm{Y})\otimes\overline{\mathrm{L}^{2}(\mathrm{X})}\) 到 \(\mathcal{L}_{2}(\mathrm{L}^{2}(\mathrm{X});\mathrm{L}^{2}(\mathrm{Y}))\) 中把 \(g\otimes f\) 映为希尔伯特--施密特映射 \(h\mapsto\langle f|h\rangle g\) 的线性映射，可扩张为一个等距同构 \(\theta_{1}\) ，即从 \(\mathrm{L}^{2}(\mathrm{Y})\widehat{\otimes}_{2}\overline{\mathrm{L}^{2}(\mathrm{X})}\) 到 \(\mathcal{L}_{2}(\mathrm{L}^{2}(\mathrm{X});\mathrm{L}^{2}(\mathrm{Y}))\) 上的等距同构（EVT, V, p. 52，定理 1）。

另外，用 \(\theta_{2}\) 表示 \(\mathrm{L}^{2}(\mathrm{X})\widehat{\otimes}_{2}\mathrm{L}^{2}(\mathrm{Y})\) 到 \(\mathrm{L}^{2}(\mathrm{Y})\widehat{\otimes}_{2}\overline{\mathrm{L}^{2}(\mathrm{X})}\) 上的等距同构，它把 \(f\otimes g\) 映为 \(g\otimes \overline{f}\) ，对所有 \(f\in \mathrm{L}^2 (\mathrm{X})\) 与每个 \(g\in \mathrm{L}^2 (\mathrm{Y})\) 成立。

线性映射 \(\theta_{3} = \theta_{1} \circ \theta_{2}\) 等同于 \(\mathrm{L}^{2}(\mathrm{X} \times \mathrm{Y})\) 到 \(\mathcal{L}_{2}(\mathrm{L}^{2}(\mathrm{X}); \mathrm{L}^{2}(\mathrm{Y}))\) 上的一个等距同构。

设 \(f \in \mathrm{L}^2(\mathrm{X})\) 且 \(g \in \mathrm{L}^2(\mathrm{Y})\) ，并设 \(\mathrm{N}\) 是 \(\mathrm{L}^2(\mathrm{X} \times \mathrm{Y})\) 的元素（与 \(f \otimes g\) 等同）。由 INT, V, p. 95, § 8, no. 3 的推论 2 及公式 (12)，我们有

\[
\langle g _ {1} \mid u _ {\mathrm{N}} (f _ {1}) \rangle = \langle \overline {{f}} \mid f _ {1} \rangle \langle g _ {1} \mid g \rangle
\]

对所有 \(f_1 \in \mathrm{L}^2(\mathrm{X})\) 与 \(g_1 \in \mathrm{L}^2(\mathrm{Y})\) 成立。映射 \(u = \theta_3(\mathrm{N})\) 就是线性映射 \(\theta_1(g \otimes \overline{f})\) ；因此它满足 \(u(h) = \langle \overline{f} | h \rangle g\) ，对所有 \(h \in \mathrm{L}^2(\mathrm{X})\) 成立。于是，对所有 \(f_1 \in \mathrm{L}^2(\mathrm{X})\) 与 \(g_1 \in \mathrm{L}^2(\mathrm{Y})\) ，可得

\[
\langle g _ {1} \mid u (h _ {1}) \rangle = \langle \overline {{f}} \mid h _ {1} \rangle \langle g _ {1} \mid g \rangle = \langle g _ {1} \mid u _ {\mathrm{N}} (h _ {1}) \rangle ,
\]

由此得 \(\theta_3(N) = u_N\) 。由于 \(\theta_3\) 与 \(N \mapsto u_N\) 都连续，我们断定 \(\theta_3(N) = u_N\) 对所有 \(N \in L^2(X \times Y)\) 成立，证明完成。

推论. --- 对每个 \(N \in L^{2}(X \times Y)\) ，线性映射 \(u_{N}\) 是希尔伯特--施密特算子，且我们有 \(\mathrm{Tr}(u_{\mathrm{N}}^{*}u_{\mathrm{N}}) = \|N\|^{2}\) 。

事实上，\(\| u\| _2 = \sqrt{\mathrm{Tr}(u^*u)}\) 对所有 \(u\in \mathcal{L}_2(\mathrm{L}^2 (\mathrm{X});\mathrm{L}^2 (\mathrm{Y}))\) 成立。

注. --- 设 \(N \in L^{2}(X \times Y)\) 。由 IV, p. 150 的推论 2，存在可数集 I、标准正交族 \((f_{i})_{i \in I}\) （\(L^{2}(X)\) 中的）与标准正交族 \((g_{i})_{i \in I}\) （\(L^{2}(Y)\) 中的），以及族 \((\alpha_{i})_{i \in I}\) （\(\mathbf{R}_{+}^{*}\) 中的族），使得

\[
u _ {\mathrm{N}} (f) = \sum_ {i \in \mathrm{I}} \alpha_ {i} \langle f _ {i} | f \rangle g _ {i}
\]

对所有 \(f \in \mathrm{L}^2(\mathrm{X})\) 成立，其中级数在 \(\mathrm{L}^2(\mathrm{Y})\) 中收敛。由 IV, p. 165 的推论 4 及上述推论，我们有

\[
\sum_ {i \in \mathrm{I}} \alpha_ {i} ^ {2} = \mathrm{Tr} (u _ {\mathrm{N}} ^ {*} u _ {\mathrm{N}}) = \| u _ {\mathrm{N}} \| _ {2} ^ {2} = \int_ {\mathrm{X} \times \mathrm{Y}} | \mathrm{N} (x, y) | ^ {2} d (\mu \otimes \nu) (x, y).
\]

另外，注意 \(h_{i,j} = \overline{f}_{i} \otimes g_{j} \in \mathrm{L}^{2}(\mathrm{X} \times \mathrm{Y})\) 对所有 \((i, j) \in \mathrm{I} \times \mathrm{I}\) 成立。于是我们有

\[
\mathrm{N} = \sum_ {i \in \mathrm{I}} \alpha_ {i} h _ {i, i}
\]

in \(\mathrm{L}^2 (\mathrm{X}\times \mathrm{Y})\) .

事实上，设 \((f_{j})_{j\in\mathrm{J}}\) 与 \((g_{k})_{k\in\mathrm{K}}\) 分别是 \(\mathrm{L}^{2}(\mathrm{X})\) 与 \(\mathrm{L}^{2}(\mathrm{Y})\) 的标准正交基，且分别是族 \((f_{i})_{i\in\mathrm{I}}\) 与 \((g_{i})_{i\in\mathrm{I}}\) 的扩张。置 \(h_{j,k}=\overline{f}_{j}\otimes g_{k}\) ，对所有 \((j,k)\in\mathrm{J}\times\mathrm{K}\) 。由引理 10，族 \((h_{j,k})_{(j,k)\in\mathrm{J}\times\mathrm{K}}\) 是 \(\mathrm{L}^{2}(\mathrm{X}\times\mathrm{Y})\) 的标准正交基。我们有 \(\langle h_{j,k}\mid\mathrm{N}\rangle=\langle g_{k}\mid u_{\mathrm{N}}(f_{j})\rangle\) ，对所有 \((j,k)\in\mathrm{J}\times\mathrm{K}\) 。若 \(j\notin I\) ，这个量为零。若 \(j\in I\) ，它等于 \(\alpha_{j}\langle g_{k}\mid g_{j}\rangle\) ，故除 k=j 外为零，而当 k=j 时 \(\langle h_{j,j}\mid N\rangle=\alpha_{j}\) 。因此

\[
\mathrm{N} = \sum_ {(j, k) \in \mathrm{J} \times \mathrm{K}} \left\langle h _ {j, k} \mid \mathrm{N} \right\rangle h _ {j, k} = \sum_ {i \in \mathrm{I}} \alpha_ {i} h _ {i, i}.
\]

